% =====================================================================
\chapter{El problema del aprendizaje supervisado}
\label{tema:1}
% =====================================================================

Este tema fija el marco común a todo el módulo: los datos se suponen extraídos de una distribución desconocida, un modelo es una función elegida dentro de una clase de hipótesis y su calidad se mide con el riesgo, la pérdida esperada sobre datos nuevos. La presentación sigue la introducción de la parte~I de los apuntes de CS229~\cite{cs229}, la teoría de la decisión estadística de Hastie, Tibshirani y Friedman~\cite[§2.4]{esl}, la formulación del capítulo~2 de las notas de MIT~6.390~\cite{mit6390} y el análisis de las pérdidas de Bishop~\cite[§1.5.5]{bishop}.

\section{Datos, hipótesis y pérdida}

\begin{definicion}[Problema supervisado]\label{def:supervisado}
Sean $\X$ un \emph{espacio de entradas} y $\Y$ un \emph{espacio de salidas}. Un \emph{problema de aprendizaje supervisado} consta de una distribución de probabilidad $\D$ sobre $\X\times\Y$, desconocida, y de un \emph{conjunto de entrenamiento} $S=\{(x_1,y_1),\dots,(x_n,y_n)\}$ formado por $n$ pares independientes con distribución $\D$. El problema se llama de \emph{regresión} si $\Y=\R$ y de \emph{clasificación} si $\Y$ es finito; en este caso sus elementos se llaman \emph{clases}.
\end{definicion}

En casi todo el módulo $\X=\R^{d}$: cada entrada es un vector de $d$ \emph{características}. Un modelo de lenguaje preentrenado resuelve un problema de clasificación en el que $x$ es el contexto, $y$ es el token siguiente y $\Y=\V$ es el vocabulario, con $K=\abs{\V}$ clases.

\begin{definicion}[Hipótesis y clase de hipótesis]
Una \emph{hipótesis} es una función $f\colon\X\to\widehat{\Y}$, donde $\widehat{\Y}$ es el espacio de predicciones ($\widehat{\Y}=\R$ en regresión; $\widehat{\Y}=\Y$ o el símplex de distribuciones sobre $\Y$ en clasificación). Una \emph{clase de hipótesis} es un conjunto $\Hc=\{f_\theta:\theta\in\Theta\}$ de hipótesis indexadas por un vector de \emph{parámetros} $\theta$.
\end{definicion}

\begin{definicion}[Pérdida, riesgo y riesgo empírico]\label{def:riesgo}
Una \emph{función de pérdida} es una aplicación $\ell\colon\widehat{\Y}\times\Y\to[0,\infty)$; $\ell(\hat y,y)$ mide el coste de predecir $\hat y$ cuando la respuesta es $y$. Las tres pérdidas de este módulo son:
\begin{enumerate}
  \item la \emph{cuadrática}, $\ell(\hat y,y)=(\hat y-y)^{2}$, para regresión;
  \item la \emph{0-1}, $\ell(\hat y,y)=\ind[\hat y\ne y]$, para clasificación con predicción de una clase;
  \item la \emph{logarítmica}, $\ell(q,y)=-\ln q_y$, para clasificación con predicción de una distribución $q$ sobre $\Y$.
\end{enumerate}
El \emph{riesgo} o \emph{error de generalización} de $f$ es $R(f)=\E_{(x,y)\sim\D}\bigl[\ell(f(x),y)\bigr]$, y su \emph{riesgo empírico} sobre $S$ es
\[
\Remp(f)=\frac{1}{n}\sum_{i=1}^{n}\ell\bigl(f(x_i),y_i\bigr).
\]
El principio de \emph{minimización del riesgo empírico} elige $\hat f=\argmin_{f\in\Hc}\Remp(f)$.
\end{definicion}

El objetivo es un riesgo pequeño, pero $\D$ es desconocida y solo $\Remp$ puede calcularse. Toda la teoría de la generalización (temas~5 y~6) estudia la diferencia entre ambos.

\section{El predictor óptimo}

Antes de restringirse a una clase $\Hc$, conviene saber cuál es la mejor predicción posible si $\D$ fuese conocida~\cite[§2.4]{esl}.

\begin{teorema}[Función de regresión]\label{teo:regresion}
Sea $\Y=\R$ con $\E[y^{2}]<\infty$, y sea $f^{*}(x)=\E[y\mid x]$. Para toda $f$ con $\E[f(x)^{2}]<\infty$, el riesgo cuadrático cumple
\[
R(f)=\E\bigl[(y-f(x))^{2}\bigr]=\E\bigl[(y-f^{*}(x))^{2}\bigr]+\E\bigl[(f^{*}(x)-f(x))^{2}\bigr].
\]
En consecuencia, $f^{*}$ minimiza el riesgo cuadrático, y el riesgo mínimo es $\E\bigl[\Var(y\mid x)\bigr]$.
\end{teorema}

\begin{proposicion}[Clasificador de Bayes]\label{prop:bayes}
Sea $\Y$ finito y $\ell$ la pérdida 0-1. El clasificador $f^{*}(x)=\argmax_{k\in\Y}P(y=k\mid x)$ minimiza el riesgo, y su riesgo es $R(f^{*})=1-\E\bigl[\max_k P(y=k\mid x)\bigr]$.
\end{proposicion}

\begin{proposicion}[La pérdida logarítmica premia la probabilidad verdadera]\label{prop:logloss}
Sea $\Y$ finito, $p(\cdot\mid x)$ la distribución condicional verdadera y $q$ cualquier distribución sobre $\Y$. Entonces
\[
\E\bigl[-\ln q_y\mid x\bigr]=H\bigl(p(\cdot\mid x)\bigr)+D_{\mathrm{KL}}\bigl(p(\cdot\mid x)\,\Vert\,q\bigr),
\]
de modo que el riesgo logarítmico se minimiza, para cada $x$, con $q=p(\cdot\mid x)$.
\end{proposicion}

\addtocounter{definicion}{1}%

\addtocounter{definicion}{1}%

\begin{table}[h]
\centering\small
\caption{Distribución conjunta $P(x,y)$ del \cref{ej:bayes}.}\label{tab:conjunta}
\begin{tabular}{@{}lccc|c@{}}
\toprule
 & $y=0$ & $y=1$ & $y=2$ & $P(x)$ \\
\midrule
$x=0$ & $0{,}2$ & $0{,}1$ & $0{,}1$ & $0{,}4$ \\
$x=1$ & $0{,}1$ & $0{,}2$ & $0{,}3$ & $0{,}6$ \\
\bottomrule
\end{tabular}
\end{table}

\section{Riesgo empírico y optimismo}

\begin{proposicion}\label{prop:insesgado}
Sea $f$ una hipótesis fija, elegida sin mirar $S$, con $\Var(\ell(f(x),y))=v<\infty$. Entonces $\E_S[\Remp(f)]=R(f)$ y $\Var_S(\Remp(f))=v/n$.
\end{proposicion}

\begin{proposicion}[Optimismo del error de entrenamiento]\label{prop:optimismo}
Sea $\Hc$ una clase de hipótesis y supóngase que $\hat f\in\argmin_{f\in\Hc}\Remp(f)$ existe. Entonces
\[
\E_S\bigl[\Remp(\hat f)\bigr]\le\inf_{f\in\Hc}R(f)\le\E_S\bigl[R(\hat f)\bigr].
\]
\end{proposicion}

La \cref{prop:optimismo} explica por qué el error de entrenamiento no sirve para evaluar un modelo: en media, subestima el riesgo del modelo ajustado. Hace falta un conjunto de datos que no haya intervenido en el ajuste.

\section{Entrenamiento, validación y prueba}

\begin{definicion}[Partición de los datos]
Los datos etiquetados disponibles se dividen al azar en tres subconjuntos disjuntos: el de \emph{entrenamiento}, con el que se ajustan los parámetros $\theta$; el de \emph{validación}, con el que se eligen los \emph{hiperparámetros} (grado de un polinomio, coeficiente de regularización, número de épocas), y el de \emph{prueba}, con el que se estima una sola vez el riesgo del modelo final~\cite[§2.8]{mit6390}.
\end{definicion}

Si se elige un hiperparámetro mirando el conjunto de prueba, ese conjunto pasa a formar parte del ajuste y, por la \cref{prop:optimismo}, la estimación final resulta optimista. La \cref{fig:particion} resume el protocolo; la validación cruzada (\cref{tema:6}) lo refina cuando hay pocos datos.

\addtocounter{figure}{1}%

\addtocounter{definicion}{1}%

